
In recent years, the rapid growth of digital technologies and the increasing availability of structured and unstructured data have transformed the way complex systems are studied.
Economic sectors that once relied mainly on descriptive statistics and expert judgment are now increasingly supported by data-driven methods.
Among these sectors, the housing market occupies a central position because it reflects household demand, credit conditions, monetary policy, and broader economic stability.

The United States housing market is influenced by multiple macroeconomic variables, including inflation, unemployment, mortgage conditions, and interest rates.
These variables do not evolve independently.
Instead, they interact in dynamic and often non-linear ways, which makes the analysis of housing trends both challenging and important.
A robust analytical framework is therefore needed in order to better understand these interactions and to build useful predictive tools.

This tutored project is part of the Data Mining and Machine Learning course.
Its main goal is to analyze and predict the evolution of housing prices in the United States using a data-driven methodology.
The project combines exploratory data analysis, predictive modeling, and dashboard development.
Beyond pure prediction, the objective is also to create a solution that remains interpretable and visually accessible.

To reach these goals, the work follows the CRISP-DM methodology, which structures the project into six phases: business understanding, data understanding, data preparation, modeling, evaluation, and deployment.
This methodological choice ensures that the work remains coherent from the definition of the problem to the interpretation of results and the presentation of the final application.

From a technical perspective, the project was implemented in Python using several major libraries.
Pandas and NumPy were used for data preparation and numerical analysis.
Scikit-learn was used for the machine learning pipeline and predictive models.
Plotly was used for interactive visualizations, and Streamlit was used to build an interactive web dashboard.
This dashboard enables users to explore the data, inspect relationships between variables, compare economic periods, and generate predictions in an intuitive way.

This report is organized into four main chapters.
The first chapter presents the general context of the project, including the problem statement, the proposed solution, and the methodology.
The second chapter reviews the theoretical background related to data mining, machine learning, regression, and dashboard visualization.
The third chapter describes the design and implementation of the proposed system, including data preprocessing, feature engineering, and the dashboard architecture.
The fourth chapter presents the results, model evaluation, discussion, and future perspectives.

Through this work, we aim to show how data mining and machine learning can be applied to a real economic problem and how interactive tools can support both analysis and decision-making.

\textbf{Integrated analytical scope.}
The work follows a complete data mining workflow: loading a housing dataset, cleaning it, analyzing correlations, building predictive models, presenting the results in a Streamlit dashboard, and enriching the analysis with data science and decision-support components. The application now includes OLAP analysis, an AI chatbot, prompt-based interaction, unsupervised learning, reinforcement-learning simulation, paper review, production-readiness validation, model governance elements, and a more complete deployment logic.

This integrated scope strengthens the nature of the project. Instead of being only a dashboard for visualization and prediction, the project becomes a small end-to-end analytical platform. It allows a user to move from raw data to business interpretation, from model training to evaluation, from segmentation to forecasting, and from technical metrics to real-world housing market comparison. The report presents these modules as part of the same coherent analytical system and explains how they improve interpretation, usability, and decision support.

The report remains organized around four chapters. Chapter 1 presents the general context and analytical scope. Chapter 2 develops the theoretical background, including OLAP, AI chatbots, unsupervised learning, reinforcement learning, and production readiness. Chapter 3 describes the implementation of the Streamlit application and its code modules. Chapter 4 discusses results, evaluation, paper comparison, business impact, and future deployment perspectives.
